\documentclass[11pt,a4paper]{article}
\usepackage[utf8]{inputenc}
\usepackage[T1]{fontenc}
\usepackage[french]{babel}
\usepackage[margin=2.4cm]{geometry}
\usepackage{amsmath,amssymb}
\usepackage{booktabs}
\usepackage{xcolor}
\usepackage[colorlinks=true,linkcolor=blue!50!black,citecolor=blue!50!black]{hyperref}
\usepackage{tcolorbox}
\tcbuselibrary{skins,breakable}
\usepackage{titlesec}
\usepackage{enumitem}
\usepackage{parskip}
\setlength{\parskip}{0.4em}

\definecolor{bleuprincip}{HTML}{1F4E78}
\definecolor{bleuclair}{HTML}{2E74B5}
\definecolor{orangealerte}{HTML}{C55A11}
\definecolor{vertok}{HTML}{548235}

\titleformat{\section}{\Large\bfseries\color{bleuprincip}}{\thesection}{0.6em}{}
\titleformat{\subsection}{\large\bfseries\color{bleuclair}}{\thesubsection}{0.6em}{}

\newtcolorbox{fichebox}[1][]{colback=vertok!6,colframe=vertok,boxrule=0.6pt,arc=2pt,
  left=7pt,right=7pt,top=5pt,bottom=5pt,breakable,#1}
\newtcolorbox{alertebox}[1][]{colback=orangealerte!7,colframe=orangealerte,boxrule=0.6pt,arc=2pt,
  left=7pt,right=7pt,top=5pt,bottom=5pt,breakable,title={\small\bfseries Approximation assumée},#1}
\newtcolorbox{noteboxbleu}[1][]{colback=bleuclair!7,colframe=bleuclair,boxrule=0.5pt,arc=2pt,
  left=6pt,right=6pt,top=4pt,bottom=4pt,breakable,title={\small\bfseries Principe},#1}

\newcommand{\cit}[1]{\textsuperscript{\hyperlink{ref#1}{[#1]}}}

\title{\vspace{-1.3cm}\textcolor{bleuprincip}{\Huge\bfseries GDPNow-Maroc}\\[0.3cm]
\textcolor{black}{\Large Reconstruction de l'IPC marocain, 2010--2026}\\[0.15cm]
\textcolor{gray}{\large Déflation des séries nominales --- méthode complète et reproductible}}
\author{}
\date{\today}

\begin{document}
\maketitle
\thispagestyle{empty}

\begin{abstract}
\noindent Ce document expose, avec toutes les formules nécessaires à sa reproduction fidèle, la construction d'une série mensuelle unique de l'Indice des Prix à la Consommation (IPC) marocain couvrant \textbf{janvier 2010 à juillet 2026} --- en raccordant deux bases officielles du HCP (base 2006 et base 2017), la première n'étant disponible qu'en moyenne annuelle sur la période qui nous intéresse. Cette série sert à déflater les indicateurs nominaux (crédit bancaire, recettes touristiques, M3, primes d'assurance) retenus dans les phases précédentes du travail.
\end{abstract}

\tableofcontents
\vspace{0.3cm}

% ============================================================================
\section{Pourquoi cette reconstruction}
% ============================================================================

\subsection{Le problème}

Plusieurs de nos indicateurs retenus sont exprimés en dirhams \textbf{courants} (nominaux) --- crédit bancaire, recettes touristiques, masse monétaire M3, primes d'assurance --- alors que la cible de chaque branche (la valeur ajoutée) est toujours en volume \textbf{réel}. Higgins\cit{1} traite ce problème explicitement à son Étape 2a :

\begin{noteboxbleu}
\textit{« Some of the monthly series are nominal and need to be deflated. »} --- Higgins (2014), Step 2a.
\end{noteboxbleu}

La formule de déflation elle-même est simple :
\begin{equation}
x_t^{\text{réel}} = \frac{x_t^{\text{nominal}}}{\text{IPC}_t / 100}
\label{eq:deflation}
\end{equation}

\subsection{Le problème pratique rencontré}

Deux séries IPC officielles existent, mais aucune seule ne couvre toute notre période d'entraînement (2010--2021) en fréquence mensuelle :
\begin{itemize}[leftmargin=1.3em]
\item \textbf{Base 2017}\cit{2} : mensuelle, mais seulement depuis janvier 2017
\item \textbf{Base 2006}\cit{3} : mensuelle en principe depuis 2007, mais les seules valeurs officiellement retrouvées pour 2010--2016 sont des \textbf{moyennes annuelles}, publiées dans des rapports régionaux du HCP
\end{itemize}

D'où la nécessité d'un raccordement (\emph{splicing}), technique que Higgins\cit{1} lui-même utilise pour ses propres séries à changement de méthodologie (SIC vers NAICS).

% ============================================================================
\section{Les données sources, avec leur origine exacte}
% ============================================================================

\subsection{Série mensuelle, base 2017 (2017--2026)}

\begin{fichebox}
115 points, janvier 2017 à juillet 2026 --- fichier fourni directement par l'utilisateur, téléchargé depuis le HCP\cit{2} (« IPC 2017 par grandes divisions, Mensuel »).
\end{fichebox}

\subsection{Moyennes annuelles, base 2006 (2010--2017)}

\begin{fichebox}
\begin{center}
\small
\begin{tabular}{lrl}
\toprule
\textbf{Année} & \textbf{IPC (base 2006)} & \textbf{Source} \\
\midrule
2010 & 108,4 & Rapport régional HCP Casablanca\cit{4} \\
2011 & 109,4 & idem \\
2012 & 110,8 & idem \\
2013 & 112,9 & idem \\
2014 & 113,4 & idem \\
2015 & 115,21 & calculé : $113{,}4 \times 1{,}016$ (HCP\cit{5} : +1,6\% en 2015) \\
2016 & 116,57 & \textbf{interpolé} (aucune source directe trouvée) \\
2017 & 117,93 & moyenne des 12 valeurs mensuelles réelles\cit{6} \\
\bottomrule
\end{tabular}
\end{center}
\end{fichebox}

\begin{alertebox}
\textbf{2016 est la seule valeur non directement sourcée} de toute la construction. Faute d'avoir trouvé une publication HCP donnant explicitement la moyenne annuelle 2016, elle est obtenue par interpolation log-linéaire entre 2015 et 2017 (formule \ref{eq:interp_2016} ci-dessous) --- une approximation isolée et documentée, pas une donnée officielle.
\end{alertebox}

\subsection{Les 12 valeurs mensuelles réelles de 2017, base 2006 --- le point de raccord}

\begin{fichebox}
Trouvées dans un rapport régional HCP (Tanger)\cit{6}, donnant le détail mensuel national complet de l'année 2017 en base 2006 :
\begin{center}
\small
\begin{tabular}{lrrrrrr}
\toprule
Jan & Fév & Mar & Avr & Mai & Jun \\
117,7 & 117,4 & 116,7 & 116,9 & 117,5 & 117,9 \\
\midrule
Jul & Aoû & Sep & Oct & Nov & Déc \\
117,3 & 117,7 & 118,7 & 118,6 & 119,1 & 119,7 \\
\bottomrule
\end{tabular}
\end{center}
\textbf{C'est cette table qui rend le raccordement particulièrement fiable} : au lieu d'un unique point de chevauchement, on dispose de 12 points communs entre les deux bases, permettant de vérifier la stabilité du coefficient de raccord plutôt que de le supposer.
\end{fichebox}

% ============================================================================
\section{Étape A --- Interpolation des moyennes annuelles (2010--2016)}
% ============================================================================

\subsection{Principe --- une variante simplifiée de Denton}

\begin{noteboxbleu}
Higgins\cit{1} utilise une interpolation de Denton pour relier une variable basse fréquence à une série liée de fréquence plus élevée (\textit{« interpolates a lower frequency variable with a related higher frequency one »}). Ici, \textbf{aucune série mensuelle liée n'est disponible} pour guider l'interpolation de 2010--2016 --- on utilise donc une variante simplifiée : une interpolation \textbf{log-linéaire} entre moyennes annuelles successives, \textbf{contrainte} pour que la moyenne des 12 valeurs mensuelles interpolées d'une année égale exactement sa moyenne annuelle officielle (c'est cette contrainte de cohérence qui constitue le vrai principe de Denton, au-delà du choix de la fonction d'interpolation elle-même).
\end{noteboxbleu}

\subsection{Formule d'interpolation}

Pour l'année $a$ (parmi 2010 à 2016) et le mois $m\in\{1,\dots,12\}$ :
\begin{equation}
\log\big(\text{IPC}_{a,m}^{\text{brut}}\big) = \log(\text{IPC}_a) + \big[\log(\text{IPC}_{a+1}) - \log(\text{IPC}_a)\big] \times \frac{m-1}{12}
\label{eq:interp_lineaire}
\end{equation}

où $\text{IPC}_a$ et $\text{IPC}_{a+1}$ sont les moyennes annuelles officielles des années $a$ et $a+1$.

\subsection{Contrainte de cohérence (la vraie logique Denton)}

Un facteur de correction $c_a$ est appliqué à chaque année pour forcer l'égalité exacte avec la moyenne officielle :
\begin{equation}
c_a = \frac{\text{IPC}_a}{\dfrac{1}{12}\sum_{m=1}^{12} \text{IPC}_{a,m}^{\text{brut}}}, \qquad
\text{IPC}_{a,m} = \text{IPC}_{a,m}^{\text{brut}} \times c_a
\label{eq:correction_denton}
\end{equation}

\begin{fichebox}
\textbf{Vérification effectuée} : après application de $c_a$, la moyenne des 12 valeurs mensuelles de chaque année reproduit exactement (à $10^{-3}$ près) la moyenne annuelle officielle, pour les 7 années 2010--2016 --- confirmé numériquement.
\end{fichebox}

\subsection{Le cas particulier de 2016}
\label{eq:interp_2016}

En l'absence de source directe, la moyenne annuelle 2016 elle-même est d'abord obtenue par interpolation log-linéaire entre 2015 et 2017 :
\begin{equation}
\log(\text{IPC}_{2016}) = \frac{\log(\text{IPC}_{2015}) + \log(\text{IPC}_{2017})}{2} \ \Rightarrow\ \text{IPC}_{2016} = 116{,}57
\end{equation}
avant d'être elle-même utilisée comme ancrage dans l'équation \ref{eq:interp_lineaire} pour interpoler ses propres 12 mois.

% ============================================================================
\section{Étape B --- Le coefficient de raccordement}
% ============================================================================

\subsection{Formule}

Sur les 12 mois de l'année de chevauchement (2017), pour chaque mois $m$ :
\begin{equation}
k_m = \frac{\text{IPC}_{2017,m}^{\text{base 2006}}}{\text{IPC}_{2017,m}^{\text{base 2017}}}
\end{equation}
Le coefficient retenu est la moyenne de ces 12 ratios :
\begin{equation}
k = \frac{1}{12}\sum_{m=1}^{12} k_m
\label{eq:coef_raccord}
\end{equation}

\subsection{Résultat}

\begin{fichebox}
$$k = 1{,}17952 \qquad \text{écart-type des 12 } k_m : 0{,}0022$$
\end{fichebox}

\begin{noteboxbleu}
L'écart-type très faible (0,22\% de la valeur du coefficient) confirme que le ratio entre les deux bases est \textbf{stable sur toute l'année de chevauchement} --- pas seulement vrai à un instant donné. C'est ce qui distingue un raccordement fiable d'un raccordement approximatif : on vérifie la stabilité plutôt que de la supposer.
\end{noteboxbleu}

% ============================================================================
\section{Étape C --- Fusion en une série unique}
% ============================================================================

\subsection{Rebasculement de la partie 2010--2016}

Chaque valeur interpolée (équation \ref{eq:correction_denton}) est ramenée à l'échelle de la base 2017 :
\begin{equation}
\text{IPC}_{t}^{\text{final}} = \frac{\text{IPC}_{t}^{\text{base 2006}}}{k} \qquad \text{pour } t < \text{janvier 2017}
\end{equation}

\subsection{La partie 2017--2026}

Conservée \textbf{sans aucune modification} --- ce sont des valeurs officielles réelles, pas des estimations :
\begin{equation}
\text{IPC}_{t}^{\text{final}} = \text{IPC}_{t}^{\text{base 2017, officiel}} \qquad \text{pour } t \geq \text{janvier 2017}
\end{equation}

\subsection{Vérification de la continuité au point de jonction}

\begin{fichebox}
\begin{center}
\begin{tabular}{lr}
\toprule
\textbf{Mois} & \textbf{IPC final} \\
\midrule
2016-10 & 99,161 \\
2016-11 & 99,257 \\
2016-12 & 99,354 \\
\textbf{2017-01} & \textbf{99,400 (valeur officielle réelle)} \\
2017-02 & 99,500 \\
\bottomrule
\end{tabular}
\end{center}
\textbf{Aucun saut visible} à la jonction --- la transition décembre 2016 → janvier 2017 (99,354 → 99,400) est cohérente avec la tendance des mois environnants, contrairement à ce qu'on observerait avec un raccordement mal calibré.
\end{fichebox}

% ============================================================================
\section{Résultat final}
% ============================================================================

\begin{fichebox}
\textbf{Série IPC complète, base 2017, mensuelle : 199 points, janvier 2010 à juillet 2026.}

\begin{center}
\small
\begin{tabular}{lll}
\toprule
\textbf{Segment} & \textbf{Nature} & \textbf{Fiabilité} \\
\midrule
2010--2016 (84 points) & Interpolé (Denton simplifié) & Approximation, contrainte aux moyennes officielles \\
2017--2026 (115 points) & Officiel, mensuel réel & Fiabilité maximale \\
\bottomrule
\end{tabular}
\end{center}
\end{fichebox}

\subsection{Application aux séries nominales}

Une fois cette série disponible, chaque indicateur nominal retenu se déflate via l'équation \ref{eq:deflation} :
\begin{equation}
x_t^{\text{réel}} = \frac{x_t^{\text{nominal}}}{\text{IPC}_t^{\text{final}} / 100}
\end{equation}

avant d'être retesté (corrélation, Student, FDR) exactement selon la même procédure que celle déjà appliquée au Niveau 2 --- les séries concernées sont le crédit bancaire (toutes branches), les recettes touristiques, la masse monétaire M3 et les primes d'assurance.

\section*{Références}
\addcontentsline{toc}{section}{Références}
\begin{enumerate}[leftmargin=1.6em, label=\textbf{[\arabic*]}]
\item \hypertarget{ref1}{} Higgins, P. (2014). \textit{GDPNow: A Model for GDP ``Nowcasting''}. Federal Reserve Bank of Atlanta, Working Paper 2014-7.
\item \hypertarget{ref2}{} Haut-Commissariat au Plan. \textit{Indice des prix à la consommation (IPC 2017) par grandes divisions, Mensuel}. Fichier fourni par l'utilisateur, source hcp.ma.
\item \hypertarget{ref3}{} Haut-Commissariat au Plan. \textit{Enquête des prix à la consommation (IPC), base 100:2006}. hcp.ma/Enquete-des-prix-a-la-consommation-IPC\_a117.html.
\item \hypertarget{ref4}{} Direction Régionale du Grand Casablanca, HCP. Note régionale sur l'IPC, tableau « Evolution annuelle de l'IPC 2006--2014 ». hcp.ma/reg-casablanca/attachment/617911/.
\item \hypertarget{ref5}{} Haut-Commissariat au Plan (2016). \textit{L'Indice des prix à la consommation (IPC) de l'année 2015}. hcp.ma/L-Indice-des-prix-a-la-consommation-IPC-de-l-annee-2015\_a1656.html.
\item \hypertarget{ref6}{} Direction Régionale de Tanger, HCP. Note régionale sur l'IPC, détail mensuel national 2017. hcp.ma/region-tanger/attachment/2066184/.
\end{enumerate}

\end{document}
